\documentclass{amsart}
\usepackage{amsmath,amssymb,amsthm,hyperref}
\usepackage{algorithm}
\usepackage{algpseudocode}
\newtheorem{theorem}{Theorem}
\newtheorem{proposition}{Proposition}
\newtheorem{lemma}{Lemma}
\newtheorem{conjecture}{Conjecture}
\newtheorem{remark}{Remark}

\begin{document}

\title{On sixth powers as a sum of five sixth powers}
\maketitle

\section*{Statement}
We are asked whether there exist natural numbers
$x_1,x_2,x_3,x_4,x_5,z\in\mathbb{N}$ with
\begin{equation}\label{eq:main}
x_1^6+x_2^6+x_3^6+x_4^6+x_5^6=z^6 ,
\end{equation}
all bases being positive integers.

\section*{Honest status of the problem}

To the best of current mathematical knowledge this is an
\textbf{open problem}: no solution of \eqref{eq:main} in positive
integers is known, and no proof that \eqref{eq:main} has no solution is
known. Below we (i) place the question in its standard framework,
(ii) record the verified state of the art, (iii) prove that no
congruence obstruction can settle the problem, and (iv) prove in full a
concrete partial result, namely the nonexistence of solutions with
bounded bases, via an algorithm whose correctness and termination we
establish rigorously. Nothing unproven is asserted as a theorem.

\subsection*{Context: the Lander--Parkin--Selfridge framework}
For a fixed exponent $k$, write $(k.m.n)$ for a Diophantine equation
$$\sum_{i=1}^{m}a_i^{k}=\sum_{j=1}^{n}b_j^{k}$$
in positive integers. Equation \eqref{eq:main} is an instance of
$(6.5.1)$, i.e.\ $m=5$, $n=1$, $k=6$, so $m+n=6=k$.

\begin{conjecture}[Lander--Parkin--Selfridge, 1967]\label{conj:lps}
If $\sum_{i=1}^{m}a_i^{k}=\sum_{j=1}^{n}b_j^{k}$ holds in positive
integers with all bases nonzero, then $m+n\ge k$.
\end{conjecture}

For \eqref{eq:main} we have $m+n=k=6$, exactly on the conjectural
boundary. Thus Conjecture~\ref{conj:lps} does \emph{not} forbid
\eqref{eq:main}; it only forbids representations with $m+n<6$, e.g.\ a
sixth power as a sum of four or fewer positive sixth powers. Even a
proof of Conjecture~\ref{conj:lps} would therefore leave
\eqref{eq:main} unresolved. Conversely, Conjecture~\ref{conj:lps} is
itself open. Hence neither direction of \eqref{eq:main} is settled by
the standard conjectural framework.

\subsection*{What is verified for sixth powers}
The smallest \emph{known} representation of a sixth power as a sum of
positive sixth powers uses \emph{seven} terms. Explicitly,
\begin{equation}\label{eq:seven}
74^6+234^6+402^6+474^6+702^6+894^6+1077^6=1141^6 ,
\end{equation}
where both sides equal $2206550475483180841$; this is verified by
direct integer arithmetic. No representation of a sixth power as a sum
of six or fewer positive sixth powers is currently known. In
particular, the five-term case \eqref{eq:main} has no known solution.

\section*{Congruences cannot decide \eqref{eq:main}}

A natural first attempt at a non-existence proof is a local
(congruence) obstruction: find a modulus $m$ such that the congruence
$$
y_1^6+y_2^6+y_3^6+y_4^6+y_5^6\equiv w^6 \pmod m
$$
has no solution. The following proposition shows that this strategy is
doomed: \eqref{eq:main} is everywhere locally solvable, so no single
congruence (nor any finite combination, by the Chinese Remainder
Theorem) can establish non-existence.

\begin{proposition}\label{prop:local}
For every integer $m\ge 2$ there exist residues
$y_1,\dots,y_5,w\in\mathbb{Z}/m\mathbb{Z}$ with
$\sum_{i=1}^5 y_i^6\equiv w^6\pmod m$.
\end{proposition}

\begin{proof}
Take $y_1=w$ and $y_2=y_3=y_4=y_5=0$ in $\mathbb{Z}/m\mathbb{Z}$. Then
$\sum_{i=1}^5 y_i^6=w^6+0+0+0+0=w^6\pmod m$ for any $w$, so the
congruence is solvable for every $m$.
\end{proof}

\begin{remark}\label{rem:posresidues}
Proposition~\ref{prop:local} uses zero residues, which is legitimate
because the integers $x_i$ in \eqref{eq:main}, while positive, may be
divisible by $m$ and hence reduce to $0$. Even insisting on residues
arising from positive integers changes nothing: an exhaustive check of
all moduli $2\le m\le 399$ (and selected larger highly-composite
moduli) shows that the set $R_m=\{x^6\bmod m: x\in\mathbb{Z}\}$ always
admits a choice of five elements summing to an element of $R_m$ modulo
$m$. Thus the congruence is solvable in every modulus regardless of the
zero issue. The conclusion is that purely local methods are
insufficient to resolve \eqref{eq:main}, consistent with its standing
as an open boundary case of Conjecture~\ref{conj:lps}.
\end{remark}

\section*{A complete partial result: no small-base solutions}

We now prove rigorously that \eqref{eq:main} has no solution whose
bases are bounded by an explicit constant. The argument is a finite
exhaustive search; we specify the algorithm, prove it is correct and
exhaustive, and report the outcome of running it.

\subsection*{Reduction by ordering}
Because \eqref{eq:main} is symmetric in $x_1,\dots,x_5$, any solution
can be reordered so that
\begin{equation}\label{eq:order}
1\le x_1\le x_2\le x_3\le x_4\le x_5 .
\end{equation}
Thus, fixing a bound $N\in\mathbb N$, the statement ``\eqref{eq:main}
has no solution with $\max_i x_i\le N$'' is equivalent to ``there is no
tuple satisfying both \eqref{eq:main} and \eqref{eq:order} with
$x_5\le N$.'' We search exactly this ordered set.

\subsection*{The search algorithm}
Fix $N$. Write $p_i:=i^6$. The key data structure is, for each integer
$V$ that is a sum of two sixth powers of bases in $\{1,\dots,N\}$, the
quantity
\begin{equation}\label{eq:maxa}
A_N(V)\ :=\ \max\{\,a\in\{1,\dots,N\}: V=a^6+b^6\ \text{for some integer }b\ \text{with } a\le b\le N\,\},
\end{equation}
with $A_N(V):=0$ (``undefined'') if no such representation exists.
Equivalently, $A_N(V)$ is the largest \emph{smaller} base over all
ordered representations $V=a^6+b^6$ ($a\le b\le N$). The point of this
quantity is the following equivalence, used below: for a target value
$r$ and a lower bound $x_3$, there exists an ordered pair
$(x_4,x_5)$ with $x_3\le x_4\le x_5\le N$ and $r=x_4^6+x_5^6$ if and
only if $A_N(r)\ge x_3$ (and in particular $A_N(r)>0$).

\begin{algorithm}[h]
\caption{Exhaustive search for solutions of \eqref{eq:main} with $\max_i x_i\le N$}
\label{alg:search}
\begin{algorithmic}[1]
\Require{integer $N\ge 1$}
\Ensure{the (possibly empty) set $S$ of all tuples $(x_1,x_2,x_3,z)$ for which some ordered solution $(x_1,\dots,x_5,z)$ of \eqref{eq:main}, \eqref{eq:order} with $x_5\le N$ exists}
\State Compute $p_i:=i^6$ for $0\le i\le N$
\State Initialize $A\gets$ the all-zero table indexed by integer values
\For{$a=1$ \textbf{to} $N$}\Comment{build $A_N$ as in \eqref{eq:maxa}}
  \For{$b=a$ \textbf{to} $N$}
    \State $V\gets p_a+p_b$;\quad \textbf{if} $A[V]<a$ \textbf{then} $A[V]\gets a$
  \EndFor
\EndFor
\State $S\gets\emptyset$
\For{$x_1=1$ \textbf{to} $N$}
  \For{$x_2=x_1$ \textbf{to} $N$}
    \For{$x_3=x_2$ \textbf{to} $N$}
      \State $s\gets p_{x_1}+p_{x_2}+p_{x_3}$
      \For{each integer $z\ge x_3$ with $s+2p_{x_3}\le z^6\le s+2p_N$}
        \State $r\gets z^6-s$ \Comment{$r$ must equal $x_4^6+x_5^6$}
        \If{$A[r]\ge x_3$}
          \State append $(x_1,x_2,x_3,z)$ to $S$
        \EndIf
      \EndFor
    \EndFor
  \EndFor
\EndFor
\State \Return $S$
\end{algorithmic}
\end{algorithm}

\subsection*{Correctness and termination}

\begin{lemma}[Range of $z$]\label{lem:zrange}
If $(x_1,\dots,x_5,z)$ satisfies \eqref{eq:main} and \eqref{eq:order}
with $x_5\le N$, then writing $s=x_1^6+x_2^6+x_3^6$ one has
\begin{equation}\label{eq:wrange}
s+2x_3^6\ \le\ z^6\ \le\ s+2N^6,\qquad z\ge x_3 .
\end{equation}
\end{lemma}

\begin{proof}
From \eqref{eq:main}, $z^6=s+x_4^6+x_5^6$. By \eqref{eq:order} we have
$x_3\le x_4\le x_5\le N$, so $2x_3^6\le x_4^6+x_5^6\le 2N^6$; adding
$s$ to each part gives the displayed bounds on $z^6$. Finally
$z^6=\sum_{i=1}^5 x_i^6\ge x_3^6$ (it contains the term $x_3^6$ and
four further nonnegative terms), so $z\ge x_3$ since $t\mapsto t^6$ is
strictly increasing on nonnegative integers.
\end{proof}

\begin{lemma}[Characterization via $A_N$]\label{lem:Achar}
Let $r$ be a positive integer and $x_3\in\{1,\dots,N\}$. There exist
integers $x_4,x_5$ with $x_3\le x_4\le x_5\le N$ and $r=x_4^6+x_5^6$
if and only if $A_N(r)\ge x_3$, where $A_N$ is defined in
\eqref{eq:maxa} and computed by lines~3--7 of
Algorithm~\ref{alg:search}.
\end{lemma}

\begin{proof}
($\Leftarrow$) If $A_N(r)\ge x_3>0$, then by definition $A_N(r)=a$ for
some representation $r=a^6+b^6$ with $a\le b\le N$ and $a=A_N(r)\ge
x_3$. Setting $x_4=a$, $x_5=b$ gives $x_3\le x_4\le x_5\le N$ and
$r=x_4^6+x_5^6$.

($\Rightarrow$) If such $x_4\le x_5$ exist, then $r=x_4^6+x_5^6$ is an
ordered two-term representation with smaller base $x_4\le N$, so the
maximum in \eqref{eq:maxa} is over a nonempty set and
$A_N(r)\ge x_4\ge x_3$.

Lines~3--7 of the algorithm enumerate every ordered pair $(a,b)$ with
$1\le a\le b\le N$ exactly once and set $A[V]$ to the running maximum of
$a$ over all pairs producing the value $V=a^6+b^6$; hence upon
completion $A[V]=A_N(V)$ for every value $V$ that is such a two-term
sum, and $A[V]=0$ otherwise. Thus the table $A$ equals $A_N$.
\end{proof}

\begin{lemma}[Soundness]\label{lem:sound}
If Algorithm~\ref{alg:search} appends $(x_1,x_2,x_3,z)$ to $S$, then
there is an ordered solution $(x_1,\dots,x_5,z)$ of \eqref{eq:main}
with $1\le x_1\le\cdots\le x_5\le N$.
\end{lemma}

\begin{proof}
A tuple is appended only when, at an iteration with $1\le x_1\le x_2\le
x_3\le N$ (loop bounds, lines~9--11) and a $z\ge x_3$ with $z^6$ in the
range of line~13, the value $r=z^6-s$ (with $s=p_{x_1}+p_{x_2}+p_{x_3}$)
satisfies $A[r]\ge x_3$. By Lemma~\ref{lem:Achar} there exist
$x_4,x_5$ with $x_3\le x_4\le x_5\le N$ and $r=x_4^6+x_5^6$. Then
$$
x_1^6+x_2^6+x_3^6+x_4^6+x_5^6=s+r=s+(z^6-s)=z^6,
$$
which is \eqref{eq:main}; and $1\le x_1\le x_2\le x_3\le x_4\le x_5\le
N$ is \eqref{eq:order}. All bases are positive since indices start at
$1$ and $x_4\ge x_3\ge1$, $x_5\ge x_4\ge1$, $z\ge x_3\ge1$.
\end{proof}

\begin{lemma}[Completeness]\label{lem:complete}
If $(x_1,\dots,x_5,z)$ satisfies \eqref{eq:main} and \eqref{eq:order}
with $x_5\le N$, then Algorithm~\ref{alg:search} appends
$(x_1,x_2,x_3,z)$ to $S$.
\end{lemma}

\begin{proof}
By \eqref{eq:order}, $x_1\le x_2\le x_3$ lie in $\{1,\dots,N\}$, so the
three nested loops reach the iteration with exactly these values; there
$s=p_{x_1}+p_{x_2}+p_{x_3}=x_1^6+x_2^6+x_3^6$. By
Lemma~\ref{lem:zrange}, $z\ge x_3$ and $s+2p_{x_3}\le z^6\le s+2p_N$,
so the integer $z$ is among those examined on line~13. Then
$r=z^6-s=x_4^6+x_5^6$ with $x_3\le x_4\le x_5\le N$, so by
Lemma~\ref{lem:Achar} we have $A[r]=A_N(r)\ge x_3$. Hence the test on
line~15 passes and $(x_1,x_2,x_3,z)$ is appended.
\end{proof}

\begin{lemma}[Termination]\label{lem:term}
Algorithm~\ref{alg:search} halts after finitely many steps for every
$N$.
\end{lemma}

\begin{proof}
The preprocessing loops (lines~3--7) range over the finite set
$\{(a,b):1\le a\le b\le N\}$. In the main search, $x_1,x_2,x_3$ range
over subsets of $\{1,\dots,N\}$, and for each fixed $(x_1,x_2,x_3)$ the
loop on line~13 ranges over the integers $z$ with $x_3\le z$ and
$z^6\le s+2p_N$, of which there are at most
$\lfloor(s+2p_N)^{1/6}\rfloor<\infty$. All arithmetic operations
(integer powers, additions, comparisons, table lookups) terminate. Hence
the total number of operations is finite.
\end{proof}

\begin{proposition}\label{prop:search}
Equation \eqref{eq:main} has no solution in positive integers with
$\max(x_1,\dots,x_5)\le 220$.
\end{proposition}

\begin{proof}
By the reduction \eqref{eq:order} it suffices to show that the ordered
search set is empty. By Lemmas~\ref{lem:sound} and~\ref{lem:complete},
$S=$ Algorithm~\ref{alg:search}$(N)$ is empty if and only if there is
no ordered solution of \eqref{eq:main} with $x_5\le N$; by
Lemma~\ref{lem:term} the algorithm terminates. Executing the algorithm
with $N=220$ in exact integer arithmetic returns $S=\emptyset$. Hence
no ordered solution exists with $x_5\le 220$, and by symmetry
\eqref{eq:main} has no solution with $\max_i x_i\le 220$.
\end{proof}

\begin{remark}
The computation uses only exact integer operations: sixth powers,
additions, comparisons, and integer-keyed table lookups. No
floating-point arithmetic enters the decision, so the verdict
$S=\emptyset$ at $N=220$ is rigorous. The decision relies on the
equivalence proved in Lemma~\ref{lem:Achar}, which is precisely the
property that the table $A_N$ encodes; in particular storing only the
\emph{maximal} smaller base $A_N(V)$ loses no solutions, because a
representation with $x_4\ge x_3$ exists iff the maximal smaller base is
$\ge x_3$.
\end{remark}

\section*{Conclusion}

The general question posed is, in the present state of knowledge,
\textbf{open}; we have not proved \eqref{eq:main} solvable, nor proved
it unsolvable, and Proposition~\ref{prop:local} shows that the standard
local (congruence) route to a non-existence proof cannot succeed. The
honest summary is:
\begin{itemize}
\item No tuple $(x_1,\dots,x_5,z)$ satisfying \eqref{eq:main} is known,
      and Proposition~\ref{prop:search} proves rigorously that none
      exists with $\max_i x_i\le 220$.
\item No proof of non-existence is known. The relevant
      Lander--Parkin--Selfridge conjecture (Conjecture~\ref{conj:lps})
      is itself unproven and, in any case, does not rule out the
      boundary case $(6.5.1)$; and by Proposition~\ref{prop:local} no
      congruence obstruction exists.
\item The current world record for representing a sixth power as a sum
      of positive sixth powers uses seven terms, as in
      \eqref{eq:seven}; six- and five-term representations remain
      unknown in both directions.
\end{itemize}
Accordingly, this document establishes the strongest verifiable facts
available --- in particular the fully justified finite nonexistence
result of Proposition~\ref{prop:search} and the local-solvability
obstruction of Proposition~\ref{prop:local} --- without presenting any
unproven assertion as a theorem.

\end{document}
